\documentclass[11pt,a4paper,oneside]{report}             % Single-side
%\documentclass[11pt,a4paper,twoside,openright]{report}  % Duplex

%\PassOptionsToPackage{chapternumber=Huordinal}{magyar.ldf}
\usepackage{t1enc}
\usepackage[utf8]{inputenc}
\usepackage{amsmath}
\usepackage{amssymb}
\usepackage{enumerate}
\usepackage[thmmarks]{ntheorem}
\usepackage{graphicx}
\usepackage{epsfig}
\usepackage{listings}
\usepackage{color}
%\usepackage{fancyhdr}
\usepackage{lastpage}
\usepackage{anysize}
\usepackage[magyar]{babel}
\usepackage{sectsty}
\usepackage{setspace}  % Ettol a tablazatok, abrak, labjegyzetek maradnak 1-es sorkozzel!
\usepackage[hang]{caption}
\usepackage{hyperref}
\usepackage{titlesec}
\usepackage{subcaption}
\usepackage{sectsty}
\usepackage{wrapfig}
\usepackage{mathtools}
\usepackage{calrsfs}
%--------------------------------------------------------------------------------------
% Main variables
%--------------------------------------------------------------------------------------
\newcommand{\vikszerzo}{Papp Dorottya}
\newcommand{\vikkonzulens}{}
\newcommand{\vikcim}{Matematikai statisztika}
\newcommand{\viktanszek}{Hálózati Rendszerek és Szolgáltatások Tanszék}
\newcommand{\vikdoktipus}{Tételsor kidolgozás}
\newcommand{\vikdepartmentr}{}

%--------------------------------------------------------------------------------------
% Page layout setup
%--------------------------------------------------------------------------------------
% we need to redefine the pagestyle plain
% another possibility is to use the body of this command without \fancypagestyle
% and use \pagestyle{fancy} but in that case the special pages
% (like the ToC, the References, and the Chapter pages)remain in plane style

\pagestyle{plain}
%\setlength{\parindent}{0pt} % áttekinthetõbb, angol nyelvû dokumentumokban jellemzõ
%\setlength{\parskip}{8pt plus 3pt minus 3pt} % áttekinthetõbb, angol nyelvû dokumentumokban jellemzõ
\setlength{\parindent}{12pt} % magyar nyelvû dokumentumokban jellemzõ
\setlength{\parskip}{0pt}    % magyar nyelvû dokumentumokban jellemzõ

\marginsize{35mm}{25mm}{15mm}{15mm} % anysize package
\setcounter{secnumdepth}{0}

\chaptertitlefont{\Large}

\setcounter{secnumdepth}{2}
\singlespacing
\frenchspacing

%--------------------------------------------------------------------------------------
%	Setup hyperref package
%--------------------------------------------------------------------------------------
\hypersetup{
    bookmarks=true,            % show bookmarks bar?
    unicode=false,             % non-Latin characters in Acrobat’s bookmarks
    pdftitle={\vikcim},        % title
    pdfauthor={\vikszerzo},    % author
    pdfsubject={\vikdoktipus}, % subject of the document
    pdfcreator={\vikszerzo},   % creator of the document
    pdfproducer={Producer},    % producer of the document
    pdfkeywords={keywords},    % list of keywords
    pdfnewwindow=true,         % links in new window
    colorlinks=true,           % false: boxed links; true: colored links
    linkcolor=black,           % color of internal links
    citecolor=black,           % color of links to bibliography
    filecolor=black,           % color of file links
    urlcolor=black             % color of external links
}

%--------------------------------------------------------------------------------------
% Set up listings
%--------------------------------------------------------------------------------------
\lstset{
	basicstyle=\scriptsize\ttfamily, % print whole listing small
	keywordstyle=\color{black}\bfseries\underbar, % underlined bold black keywords
	identifierstyle=, 					% nothing happens
	commentstyle=\color{white}, % white comments
	stringstyle=\scriptsize\sffamily, 			% typewriter type for strings
	showstringspaces=false,     % no special string spaces
	aboveskip=0pt,
	belowskip=0pt,
	columns=fixed,
	backgroundcolor=\color{lightgray},
} 		
\def\lstlistingname{lista}	

%--------------------------------------------------------------------------------------
%	Some new commands and declarations
%--------------------------------------------------------------------------------------
\newcommand{\code}[1]{{\upshape\ttfamily\scriptsize\indent #1}}

% define references
\newcommand{\figref}[1]{\ref{fig:#1}.}
\renewcommand{\eqref}[1]{(\ref{eq:#1})}
\newcommand{\listref}[1]{\ref{listing:#1}.}
\newcommand{\sectref}[1]{\ref{sect:#1}}
\newcommand{\tabref}[1]{\ref{tab:#1}.}

\DeclareMathOperator*{\argmax}{arg\,max}
%\DeclareMathOperator*[1]{\floor}{arg\,max}
\DeclareMathOperator{\sign}{sgn}
\DeclareMathOperator{\rot}{rot}
\definecolor{lightgray}{rgb}{0.95,0.95,0.95}

\author{\vikszerzo}
\title{\viktitle}
\includeonly{
	titlepage,%
	alapfogalmak,%
	becslesek
}
%--------------------------------------------------------------------------------------
%	Setup captions
%--------------------------------------------------------------------------------------
\captionsetup[figure]{
%labelsep=none,
%font={footnotesize,it},
%justification=justified,
width=.75\textwidth,
aboveskip=10pt}

\renewcommand{\captionlabelfont}{\small\bf}
\renewcommand{\captionfont}{\footnotesize\it}

%--------------------------------------------------------------------------------------
% Table of contents and the main text
%--------------------------------------------------------------------------------------
\begin{document}
\singlespacing

\pagenumbering{arabic}

\linespread{1.0}
\setlength{\parindent}{0.01em}
\setlength{\parskip}{0.01em}
\titlespacing*{\section}
{0pt}{1.0em plus .1em minus .2em}{.5em plus .2em}
\titlespacing*{\subsection}
{0pt}{1.0ex plus .1em minus .2em}{.5em plus .2em}

%--------------------------------------------------------------------------------------
%	The title page
%--------------------------------------------------------------------------------------
\begin{titlepage}
\begin{center}
\includegraphics[width=60mm,keepaspectratio]{figures/BMElogo.png}\\
\vspace{0.3cm}
\textbf{Budapesti Műszaki és Gazdaságtudományi Egyetem}\\
\textmd{Villamosmérnöki és Informatikai Kar}\\
\textmd{\viktanszek}\\[5cm]

\vspace{0.4cm}
{\huge \bfseries \vikcim}\\[0.8cm]
\vspace{0.5cm}
\textsc{\Large \vikdoktipus}\\[4cm]

\begin{tabular}{cc}
 \makebox[14cm]{\emph{Készítette}}\\
 \makebox[14cm]{\vikszerzo}
\end{tabular}

\vfill
{\large \today}
\end{center}
\end{titlepage}



\tableofcontents\vfill
\chapter{Alapfogalmak}

A \emph{sokaság}, vagy populáció, a vizsgálat tárgyát képező nagyszámú, de véges elemszámú egyedek halmaza. A statisztika, mint tudomány célja a halmaz egészének kevés adattal történő tömör jellemzése és az egyedek leírására a bevezetett változók közötti kapcsolatok leírása. Nincs lehetőség (erőforrás), hogy a populáció minden egyes eleméről adatokat szerezzünk be. Példa sokaságok: ország állampolgárai, gyártott termékek. A statisztikai elemzés tárgya lehet egy \emph{véletlen kísérlet} is, ami időben változatlan körülmények között elvileg akárhányszor lejátszódhat, pl. terméshozam, lottóhúzás, véletlen egyed kiválasztása.

A statisztikai minta \emph{realizáltja} a populáció egy kis elemszámú részhalmazára vonatkozó megfigyelések adatai. Például: hőmérséklet adatok, felmérésbe bevont nézők. A mintának reprezentatívnak kell lennie, vagyis tükröznie kell a populáció tulajdonságait. A populáció minden egyedének ugyanakkora esélyt kell biztosítani a mintába kerüléshez, a minta elemszámának pedig elég nagynak kell lennie ahhoz, hogy a következtetések átvihetők legyenek a populációra is. A mintavételezésnek 3 nagy kategóriája:
\begin{itemize}
\item Rétegzett mintavételezés: populációt adott szempontok szerint csoportokba osztjuk és a csoportok arányait a mintában is megtartjuk
\item Véletlen mintavételezés: mintába kerülő egyedeket sorsolással választjuk ki
\item Cenzus: népszámlálás
\end{itemize}
A mintát egy adatmátrixba írhatjuk fel, aminek $n$ sora és $p$ oszlopa van. Minden sor megfelel a minta egy elemének, minden oszlop pedig egy \emph{változónak}, a populáció egy mérhető jellemzőjének. 

A \emph{statisztika} a minta realizáció adataiból adott képlettel számolt adat. Pl. átlag, medián, gyakoriság, stb.

\section{Alapmodell}

A matematikai statisztika alapmodellje 4 elemből áll:
\begin{itemize}
\item $\mathcal{K}$, a véletlen kísérlet,
\item $\Omega$, a lehetséges kimenetelek halmaza,
\item $\mathcal{A}$, a megfigyelhető események halmaza, és 
\item $\mathcal{P}$, a lehetséges valószínűségi mértékek halmaza.
\end{itemize}
Az elemzés célja, hogy $\mathcal{P}$-ből kiválasszuk a tényleges valószínűséget, de legalább is egy jó helyettesítő egyedet.

A matematikai alapmodellben a változó $X: \Omega \rightarrow \mathbb{R}$ egy valószínűségi változó, melynek minden $\mathbf{P} \in \mathcal{P}$ esetén megadható az eloszlásfüggvénye: $F_x(t)=\mathbf{P}(X<t)$. Ekkor a feladat a következő halmazból kiválasztani a vlóságot legjobban leíró eloszlásfüggvényt:

$\mathcal{F}=\{F_x(t):F_x(t)=\mathbf{P}(X<t): \forall \mathbf{P} \in \mathcal{F}\}$

A statisztikai minta ez alapján az $X$ valószínűségi változóval azonos eloszlású, egymással teljesen független $X_1, X_2, ... X_n$ valószínűségi változók együttese. Egy mintavételkor tulajdonképpem megfigyeljük a $\mathcal{K}$ véletlen kísérletet, aza megállapíTjük, hogy melyik $\omega \in \Omega$ kimenetel realizálódott. Az $X_1(\omega)=x_1, ..., X_n(\omega)=x_n$ szám n-est nevezzük a minta realizációjának.

A statisztika egy valószínűségi változó, aminek eloszlásfüggvényét a minta eloszlásfüggvényéből lehet kiszámolni: $T_n=t_n(X_1,X_2,...,X_n)$. A statisztika számolt értéke az, amikor $t_n()$ argumentumába a mintarealizáció értékeit helyettesítjük.

Adatcentrum statisztikái:
\begin{itemize}
\item átlag: $\displaystyle\frac{1}{n} \sum_{i=1}^n x_i$
\item medián: $ m_n =
  \begin{cases}
    x_\frac{n+1}{2}       & \quad \text{ha } n \text{ páratlan}\\
    \frac{x_\frac{n}{2}^*+x_\frac{n+1}{2}^*}{2}  & \quad \text{ha } n \text{ páros}\\
  \end{cases}
$
\item módusz: leggyakrabban előforduló érték a mintában
\end{itemize}

Szóródást jellemző statisztikák:
\begin{itemize}
\item standard szórás: $\sqrt{\frac{1}{n-1}\sum_{i=1}^n(x_i - \bar{x})^2}$, ahol $\bar{x}$ az átlag
\item variáció: $\frac{1}{n-1}\sum_{i=1}^n(x_i - \bar{x})^2$
\item terjedelem: $x_n^*-x_1^*$ (legkisebb és legnagyobb elem közötti különbség)
\end{itemize}

Eloszlást jellemző statisztikák:
\begin{itemize}
\item ferdeség: $s=\frac{\frac{1}{n}\sum_{i=1}^n(x_i-\bar{x})^3}{\Big(\sqrt{\frac{1}{n-1}\sum_{i=1}^n(x_i - \bar{x})^2}\Big)^3}$, ha $s>0$, akkor jobbra ferdül, ha $s<0$, akkor balra ferdül
\item lapultság: $k=\frac{\frac{1}{n}\sum_{i=1}^n(x_i-\bar{x})^4}{\Big(\sqrt{\frac{1}{n-1}\sum_{i=1}^n(x_i - \bar{x})^2}\Big)^4}-3$, ha $k>0$, akkor csúcsos, ha $k<0$, akkor lapos.
\end{itemize}

A mintaelemeket szokás rendezni is: $X_k^* = ord_k(X_1,X_2,...,X_n), k=1,2,...,n$\footnote{$ord()$ kiválasztja a $k$-dik legnagyobb elemet az argumentumok közül}. Ezeket rendezett mintaelem-statisztikáknak hívjuk. A minta eloszlásfüggvénye és a rendezett minta eloszlásfüggvénye között az alábbi összefüggés áll fenn:

$F_k(x) = \mathbf{P}(X_k^* < x) = \sum_{i=k}^n \binom{n}{i} [F(x)]^i [1-F(x)]^{n-i}$

A rendezett mintaelem-statisztikák közül különlegesen fontos az empirikus eloszlásfüggvény:

$F_n(x) = 
  \begin{cases}
    0       & \quad \text{ha } x \leq X_1^* \\
    \frac{k}{n}  & \quad \text{ha } X_k^* < x \leq X_{k+1}^*\\
    1		& \quad \text{ha } x>X_n^* \end{cases}
  $
  
Fontosságát a matematikai statisztika alaptétele adja:

$\mathbf{P}\Big( \displaystyle\lim_{n \to \inf
} sup_{x\in\mathbb{R}} |F_n(x)-F(x)| = 0 \Big) = 1$, vagyis az empirikus eloszlásfüggvény 1 valószínűséggel, egyenletesen konvergál az elolszlásfüggvényhez.

\section{Fontos eloszlások}

\begin{itemize}
\item $\chi^2$ eloszlás
\item Student-eloszlás
\item F-eloszlás
\end{itemize}
\chapter{Statisztikai becslések}

Előfordul, hogy a minta eloszlásfüggvénye egy képlet, amelyet a $\vartheta$ paraméter konkretizál. Ha ismerjük az értékét, meg tudjuk pontosan adni az eloszlásfüggvényt: 

$\mathcal{F} = \{ F_x(t,\vartheta):\vartheta\in \Theta \}$

Egy adott statisztikai minta segítségével a $\vartheta$ paraméter megbecslése a célunk. A paraméter becsléséhez valamilyen alkalmas $T_n$ statisztikát használunk, egy ismeretlen számot egy valószínűségi változóval becsülünk. Ekkor persze felvetődik a kérdés: mikor jó egy ilyen becslés?
\begin{itemize}
\item Torzítatlanság: $\mathbf{E}T_n=\vartheta$, vagyis a becslő statisztika pont a becsülendő paraméterérték körül fogja felvenni az értékeit. Indok: egy valószínűségi változó az összes szám körül épp a várható érték körül ingadozik a legkisebb mértékben
	\begin{itemize}
	\item Aszimptotikus torzítatlanság: a torzítatlansági feltétel csak $n \to \infty$ esetében igaz: $\lim_{n \to \infty} \mathbf{E}T_n = \vartheta$
	\item A Cramer-Rao egyenlőtlenség elvi alsó korlátot ad a torzítatlan becslések szórásnégyzeteire. Ha egy statisztikára belátjuk, hogy a szórásnégyzete éppen az alsó korláttal egyenlő, akkor az biztosan hatásos (amiből egyetlen egy van)
	\end{itemize}
\item Konzisztencia: $\forall \epsilon >0 \text{ szám esetén teljesül, hogy } \lim_{n \to \infty}\mathbf{P}(|T_n-\vartheta|<\epsilon) = 0$, vagyis garancia van arra, hogy a minta elemszám növekedtével növekszik a becslés pontosságának valószínűsége
	\begin{itemize}
	\item Erős konzisztencia: $\mathbf{E}
	T_n = \vartheta \text{ és } \lim_{n \to \infty} \mathbf{D}^2T_n = 0$, vagyis azok a torzítatlan becslések, ahol az elemszám növekedtével a szórásnégyzet (variancia) 0-hoz tart. Az erősen konzisztenc statisztikai becslések egyben konzisztensek is
	\end{itemize}
\item Hatásosság: két torzítatlan becslés közül a kisebb varianciájú a jobb, hiszen kisebb mértékben ingadozik a paraméter körül, vagyis kevesebb megfigyeléssel is jó becslés kapható. Egy torzítatlan becslés akkor lesz \emph{hatásos}, ha varianciája minden más torzítatlan becslés varianciájánál kisebb. Hatásos becslésből egyetlen egy létezik csak, ezt érdemes megkeresni egy adott paraméter-becslési problémához
\item Elégségesség: egymaga képes helyettesíteni a mintát, a paraméterre vonatkozóan minden információt magába sűrít
	\begin{itemize}
	\item Rao-Blackwell-Kolmogorov tétel: ha létezik hatásos (legjobb torzítatlan) becslés, akkor elég azt az elégséges becslés függvényei között keresni
	\item Ellenőrzini a Neymann-Fisher faktorizációs tétellel lehet
	\end{itemize}
\end{itemize}

Jó tudni:
\begin{itemize}
\item Az átlagstatisztika a várható értéknek, mint paraméternek, a torzítatlan becslése (ha létezik a minta szórása, akkor erősen konzisztens), normális, exponenciális és Poisson esetben hatásos
\item A tapasztalati szórásnégyzet a minta varianciájának aszimptotikusan torzítatlan becslése (ha létezik a minta negyedik momentuma, akkor konzisztens is)
\item A korrigált empirikus szórásnégyzet statisztik a variancia torzítatlan becslése (ha létezik a minta negyedik momentuma, akkor erősen konzisztens), normális esetben hatásos
\end{itemize}

\section{Maximum likelihood becslés}

A módszer alapgondolata, hogy ha egy kísérletnél több esemény is bekövetkezhet, akkor legtöbbször a legnagyobb valószínűségű esetményt fogjuk megfigyelni. A mintavételezés során kaptunk egy realizációt, errőr a realizációról pedig az előzőek miatt feltételezzük, hogy azért ezt kaptuk, mert ennek volt a legnagyobb bekövetkezési valószínűségee. Vagyis, az összes lehetséges $\vartheta$ paraméter (amitől függ a minta eloszlásfüggvénye) közül vegyük azt, amelynél a kapott realizáció bekövetkezése a maximális. Általános feltételed mellett a becslés konzisztens, aszimptotikusan normális eloszlású és ha van elégséges statisztikai, akkor éppen azt adja meg.

Számítása diszkrét esetben:

$
L(\mathbf{x}, \vartheta) = P_\vartheta(X_1=x_1, X_2 = x_2, ..., X_n=x_n) = \sum_{i=1}^nP_\vartheta(X_i=x_i)$, a minta együttes eloszlása. A paraméter maximum likelhood becslése az a $\tau_n(X_1,X_2, ..., X_n)$ statisztika, melyre $L(x, \tau_n(\mathbf{x})) = \text{max}_\vartheta L(\mathbf{x}, \vartheta)
$

Számítása folytonos esetben:

$
L(\mathbf{x}, \vartheta) = P_\vartheta(X_1=x_1, X_2 = x_2, ..., X_n=x_n) = \prod_{i=1}^nf_\vartheta(x_i)$, a minta együttes eloszlása. A paraméter maximum likelhood becslése az a $\tau_n(X_1,X_2, ..., X_n)$ statisztika, melyre $L(x, \tau_n(\mathbf{x})) = \text{max}_\vartheta L(\mathbf{x}, \vartheta)
$

A maximum meghatározásához $L(\mathbf{x}, \vartheta)$ $\vartheta$ szerinti első deriváltjának zérushelyeit keressük. A zérushelyek közül az a maximum, ahol a második derivált negatív. Megjegyzés: ha $L(\mathbf{x}, \vartheta)$ deriválása túl nehéz, érdemes a logaritmusát venni (log-likehood függvény, $l(\mathbf{x}, \vartheta)$) és azzal számolni. Mivel a logaritmus függvény szigorúan monoton nő, ezért nem torzítja a zérushelyeket.

Többparaméteres esetben paraméterenként vizsgáljuk az első deriváltakat és a Hesse-mátrix alapján döntünk a maximum helyről. Maximum esetén a mátrix diagonális és az első diagonál elem negatív.

\section{A momentumok módszere}

Legyen adott a valószínűségi mértékek egy tere és az $X_1, X_2, ..., X_n$ statisztikai minta. Tegyük fel, hogy létezik az első $k$ momentum: $m_j = \mathbf{E}_\mathbf{\vartheta}X_i^j = g_j(\vartheta)$, vagyis a momentumok a paraméter valamilyen $g()$ függvényeként állnak elő. Tegyük fel, hogy $\exists g_j^{-1}(m_1, m_2, ..., m_k) = \vartheta_j$. Tekintsük az $\hat{m}_j=\frac{1}{n}\sum_{i=1}^nX_i^j$ empirikus momentum statisztikákat. Ekkor  a $\vartheta_j$ paraméteres momentumos becslései a $m_j = g_j^{-1}(\hat{m}_1, \hat{m}_2, ...,\hat{m}_k)$

\section{Intervallumbecslések}

Eddig az ismeretlen paramétervektort a minta egy függvényével, azaz egyetlen statisztikával próbáltuk meg közelíteni. Konkrét realizációnál tehát, a paramétertér egy pontját egy másik ponttal becsüljük (\emph{pontbecslés}).

Folytonos eloszlásoknál azonban annak valószínősége, hogy a valószínőségi változó az értékkészletének éppen egy tetszőlegesen kiválasztott pontját fogja felvenni, nulla. Tehát folytonos esetben nulla annak valószínősége, hogy éppen a paramétert találtuk el a becsléssel. Az intervallumbecsléseknél a mintából készített tartományokat definiálunk, amely tartományok nagy valószínőséggel lefedik a kérdéses paraméterpontot.

Legyen adott a valószínésgi mértékeke egy tere és az $X_1, X_2, ..., X_n$ statisztikai minta. Legyen $0<\epsilon<1$ rögzített. A $\vartheta$ paraméterhez megadhatunk egy legalább $1-\epsilon$ szignifikanciaszintű konfidencia-intervallumot, ha $t_1(X_1, X_2, ..., X_n)$ és $t_2(X_1, X_2, ..., X_n)$ olyan statisztikák, hogy:

$\mathbf{P}_\vartheta(t_1(X_1, X_2, ..., X_n) \leq \vartheta \leq t_2(X_1, X_2, ..., X_n)) \geq 1-\epsilon$ mindig igaz.

%\listoffigures\addcontentsline{toc}{chapter}{Ábrák jegyzéke}
%\listoftables\addcontentsline{toc}{chapter}{Táblázatok jegyzéke}

\bibliography{mybib}
\addcontentsline{toc}{chapter}{Irodalomjegyzék}
\bibliographystyle{plain}

\include{appendices}

\label{page:last}
\end{document}
